% Typeset form of COVER_LETTER.md, kept in sync by hand: the markdown is the prose record, this file is what gets
% uploaded as a PDF. Figures quoted here are the same ones the manuscript's gate checks.
\documentclass[11pt]{article}
\usepackage[a4paper,margin=2.5cm]{geometry}
\usepackage[T1]{fontenc}
\usepackage[utf8]{inputenc}
\usepackage{newtxtext}
\usepackage{parskip}
\usepackage{url}
\pagestyle{empty}
\begin{document}

\begin{center}{\Large\bfseries Cover letter --- JVCIR submission}\end{center}

\textbf{Manuscript:} CARE-GS: coverage-aware rate reduction for 3D Gaussian Splatting assets\par
\textbf{Author:} Tianji Su, Ningbo University\par
\textbf{Type:} Regular paper\par
\textbf{Date:} 2026-11-18\par

\medskip\hrule\medskip

Dear Editor,

Please consider the enclosed manuscript for publication in the \emph{Journal of Visual Communication and Image
Representation} as a regular paper.

\textbf{What the paper does.} A 3D Gaussian Splatting asset is now a file that changes hands: it is trained once,
shipped, and then re-compressed by whoever receives it. The sender and the receiver share only fidelity numbers ---
PSNR, SSIM, LPIPS --- and those numbers do not say whether the per-pixel quality guarantee that travelled with the
asset still holds. The manuscript measures what a recipient's edit does to that guarantee, and turns the measurement
into a method: a container that makes attribute quantisation worth spending ($17.35\times$ at $-0.04$~dB), a
coverage-aware planner that chooses the operating point against a stated tolerance, and an audit that the recipient
can run on the calibration views the asset already carries ($56/56$ decisions, no false acceptance, no new labels).

\textbf{Where it sits in the journal.} The work is visual data reduction and compression for a 3D representation, and
it is measured in the vocabulary this journal's compression papers use: a bitstream composition table, BD-rate for the
pruning arms ($-45.06\%$), and a decode-time and memory budget. The journal's 3DGS line has so far addressed how such
an asset is \emph{trained and optimised}; this manuscript addresses what happens to it after it leaves the training
machine, and the manuscript states that relationship explicitly, with three of the journal's recent 3DGS papers cited
in \S2.1.

The comparison literature for this particular step lives outside the venue (anchor-based and progressively quantised
3DGS coders), so rather than leave that comparison to the reader, the gap table computes it from the published numbers
and reports where this container loses: across four published coders at their own operating points the gap is
$0.96$--$3.62$~dB on \emph{bicycle} ($2.97$~dB behind HAC at $25.67$~MB, $3.62$~dB behind PCGS at $18.40$~MB) and
$6.25$~dB behind HAC on \emph{stump}. The claim is accordingly narrow --- not that the rate--distortion frontier is
beaten, but that rate reduction can be certified, and what certification costs. The Limitations section states the
remaining boundaries in the same terms (no subjective study, a per-pixel rather than view-level guarantee, and no
closure of the fidelity gap).

\textbf{What is new, in three sentences.}
\begin{enumerate}\setlength{\itemsep}{1pt}
  \item The damage an edit does to a shipped coverage certificate is decided by the \emph{action family}, not by its
        strength: inside a one-point tolerance, quantisation reaches $8.97$--$15.78\times$ on six scene-checkpoints
        while pruning has no admissible configuration at all.
  \item A shared rank-16 colour basis, smallest-three quaternions and per-stream alphabets take a $1447$~MB asset to
        $83.41$~MB at $24.08$~dB, with every step attributed in an ablation table.
  \item The planner is admissible only because the audit exists: a leave-one-asset-out prior reaches
        $0.91$--$1.06\times$ the oracle rate but violates the tolerance on two to three of six assets, whereas the
        audit on the shipped calibration views decides all $56$ measured cells with zero false accepts.
\end{enumerate}

\textbf{Data and code.} Under the journal's Option C, the code, the per-view residual dumps, the derived tables and the
reproduction scripts are archived at Zenodo and maintained at
\url{https://github.com/11qishiyi/care-gs} (release \texttt{v1.0.0}), with a permanent snapshot at Software Heritage.
Every number in the tables is regenerated from those artefacts by a script in the archive, and a verification gate in
the repository re-checks the manuscript against them.

\textbf{Suggested reviewer profiles.} Because the topic sits between two communities, I would suggest: one reviewer
from the journal's 3DGS cluster (dynamic or feed-forward Gaussian splatting), one from the journal's point-cloud codec
cluster (G-PCC / V-PCC rate--distortion or complexity work), and one from neural-field compression outside the venue
(anchor-based or vector-quantised 3DGS coding). I have no conflicts to declare with respect to these areas.

\textbf{Declarations.} The manuscript is original, has not been published elsewhere, and is not under consideration by
another journal. The author declares no competing interests. A generative-AI declaration is included in the
manuscript, as required.

Thank you for your consideration.

\medskip
Yours sincerely,\par
Tianji Su\par
Ningbo University, Ningbo, China\par
\url{71736850@qq.com}

\end{document}
